

\section{Preliminaries} \label{sec:prelim}


\subsection{Proof systems}
All our rings are commutative and unital. {\em Polynomial calculus} (PC) is a proof system to prove that a given system of (multivariate) polynomial equations ${\cal P}$ over a field $\mathbb F$ of the form $p=0$, has no solution over the algebraic closure $\overline{\F}$ (i.e. the system is unsolvable). We usually shorten the polynomial equation $p=0$ to just $p$.
The derivation rules of the system are the following one:
$$
\frac{p}{xp} \; \mbox{(multiplication)}, \quad \frac{p\quad q}{ap+bq} \; \mbox{(linear combination)} 
$$
where $x$ is any formal variable, $a,b \in \mathbb{F}$ and $p,q$ are polynomials over $\mathbb F$.  

When refuting Boolean systems of equations it is common to include the Boolean axioms $x_i^2 - x_i$. Because we do \emph{not} always include these (esp. for \textsc{TI}) we are explicit about our use of these, but do not assume they are built into the proof system---that is, if we are assuming them as axioms, we say so.

A {\em PC derivation (or proof)}  of a  polynomial $q$ from a set of polynomials ${\cal P}$ is a sequence of polynomial equations $p_1,\ldots, p_{m}$ ending with the polynomial $q$ (so $p_m$ is $q$) and where each $p_i$, $i \in [m]$, is either an {\em axiom} $p$ for $p \in  {\cal P}$, or is obtained from previous equations in the refutation by multiplication or linear combination.  We denote this by writing  ${\cal P}\vdash q$.  Observe that if $p$  is derivable in  PC and $q$ is a polynomial then, by repeated applications of multiplication and linear combination rules, we can derive
$pq$. We often use this generalization of the multiplication in our proofs without mention.  

A {\em PC refutation}  is just a PC proof  of the polynomial $1$.   The {\em degree} of a PC derivation is the maximal degree of a polynomial used in the proof. The {\em size} of a polynomial $p$ is the number of  terms in $p$.  The {\em size} a PC derivation $p_1,\ldots, p_m$  is the sum of the sizes of the polynomials $p_1,\ldots,  p_m$. 


\medskip

For our upper bound in Theorem~\ref{thm:upper}, we also consider another algebraic proof system, known as {\em Nullstellensatz} (NS), to certify unsolvability of sets of polynomial equations. Nullstellensatz  is defined in a static form as follows:  a  refutation of a list  ${\cal P}=(p_1,\ldots,p_m)$ of polynomial equations over variables $x_1\ldots,x_n$ 
is given by the list of polynomials  ${\cal Q}=(q_1,\ldots q_m)$  such that
$$
\sum_{i \in [m]} p_iq_i =1
$$

The {\em degree} of a NS refutation is the maximal degree of a polynomial in ${\cal P} \cup {\cal Q}$. The size of $NS$ proof is the sum of the number of monomials appearing in the polynomials
$q_1,\ldots, q_m$.

\medskip

{\em Sum-of-Squares} (SOS) is a static proof system for certifying the unsolvability of systems of polynomial equations and polynomial inequalities, where polynomials are usually over the ring $\mathbb R[x_1\ldots, x_n]$.    

A polynomial $p$ is a {\em sum-of-squares} polynomial if it is in the form $p=\sum_{i} r_i^2$ and the $r_i$'s  are polynomials as well.  Given a system made by a set of polynomial equations  
 ${\cal P} =\{p_1=0,\ldots p_m=0\}$ and a set ${\cal Q}=\{q_1\geq 0,\ldots q_k\geq 0\}$ of polynomial inequalities, a {\em sum-of-squares} proof of the polynomial inequality $p\geq 0 $ from ${\cal P} \cup {\cal Q}$  is  given by the  formal identity
 $$
 p= s_0+\sum_{i\in [k]}s_iq_i + \sum_{j \in [m]}t_jp_j
 $$
where $s_0,s_1,\ldots, s_k$ are sum-of-squares polynomials, while $t_1,\ldots,t_m$ are arbitrary polynomials.  When the system  ${\cal P} \cup {\cal Q}$  is unsatisfiable, a {\em refutation} of ${\cal P} \cup {\cal Q}$ is a proof of the inequality $-1\geq 0$, that is for  $p$ the constant polynomial $-1$. The {\em degree} of the proof is 
the $\max \{\deg(p),\deg(s_0), \deg(s_i)+\deg(q_i), \deg(t_j)+\deg(p_j) | i \in [k], j \in[m]\}$.   


\medskip 

\begin{definition}[PC reduction between systems of polynomials, cf. {\cite[Sec.~3]{DBLP:journals/jcss/BussGIP01}}] \label{def:reduction}
Let $P(x_1,\ldots,x_n)$ and $Q(y_1,\ldots, y_m)$ be two sets of polynomials over a field $\mathbb F$.  
$P$ is \emph{$(d_1,d_2)$-reducible} to $Q$ if:
\begin{enumerate}
\item For each $i \in [m]$ there is a polynomial $r_i(\mathbf{x})$ of degree at most $d_1$  (which we think of as defining $y_i$ in terms of the $\mathbf x$ variables);
\item There exists a degree $d_2$ PC derivation of $Q(r_1(\mathbf x),\ldots, r_m(\mathbf x) )$ from polynomials $P(\mathbf x)$.
\end{enumerate}
\end{definition}

\begin{lemma}[{\cite[Lem.~1]{DBLP:journals/jcss/BussGIP01}}]
\label{lem:PCred}
If $P(\mathbf{x})$ is $(d_1, d_2)$-reducible to $Q(\mathbf{y})$ and there is
a degree $d$ PC refutation of $Q(\mathbf{y})$, then there is a degree $\max(d_2, d_1 d)$ 
refutation of $P(\mathbf x)$.
\end{lemma}

In their paper, they typically only applied this to systems of equations which were known to be unsatisfiable (such as PHP and Tseitin tautologies), whereas in our paper we have several situations we want to combine the above notion together with the usual notion of many-one reduction. We encapsulate this in the following definition. We say a decision problem $\Pi$ is a \emph{polynomial solvability problem} over a field $\F$ if all valid instances of the problem are systems of polynomial equations over $\F$, and the problem is to decide whether such a system of equations has solutions over the algebraic closure $\overline{\F}$. Thus, the difference between multiple polynomial solvability problems is just \emph{which} systems of equations are valid inputs.

\begin{definition}[PC many-one reduction] \label{def:reduction2}
Let $\Pi_1, \Pi_2$ be two polynomial solvability problems over a field $\F$. We say that $\Pi_1$ \emph{$(d_1,d_2)$-many-one reduces} to $\Pi_2$ if there is a polynomial-time many-one reduction $\rho$ from $\Pi_1$ to $\Pi_2$, such that for all unsatisfiable instances $\mathcal{F}$ of $\Pi_1$, $\mathcal{F}$ $(d_1, d_2)$-reduces to $\rho(\mathcal{F})$. When this occurs with $d_1,d_2 = O(1)$, we write \[\Pi_1 \leq_m^{PC} \Pi_2.\]
\end{definition}



\subsection{Linear algebra and tensors}
Given three vector spaces $U,V,W$ over a field $\F$, a 3-tensor is an element of the vector space $U \otimes V \otimes W$, whose dimension is $(\dim U)(\dim V)(\dim W)$. If $e_i$ is the $i$-th standard basis vector, then a basis for $U \otimes V \otimes W$ is given by the vectors $\{e_i \otimes e_j \otimes e_k\}$. One may also interpret the symbol $\otimes$ more concretely as the Kronecker product, in which $e_i \otimes e_j \otimes e_k$ represents a 3-way array whose only nonzero entry is in the $(i,j,k)$ position. The vector space of such 3-way arrays (with coordinate-wise addition) is isomorphic to $U \otimes V \otimes W$.

The rank of a tensor $T \in U \otimes V \otimes W$ is the minimum $r$ such that $T = \sum_{i=1}^r u_i \otimes v_i \otimes w_i$ for some vectors $u_i, v_i, w_i$. 

Two $n \times m \times p$ 3-tensors $T,T' \in U \otimes V \otimes W$ are isomorphic if there exist matrices $X \in \GL(U), Y \in \GL(V), Z \in \GL(W)$ such that $(X,Y,Z) \cdot T = T'$, where the latter is shorthand for (\ref{eq:TI}).
If we treat $T,T'$ as given non-isomorphic tensors, then we may treat (\ref{eq:TI}) as a system of equations in the $n^2 + m^2 + p^2$ variables $X_{ii'}, Y_{jj'}, Z_{kk'}$. To enforce that these variable matrices are invertible, we furthermore introduce three additional sets of variables $X', Y', Z'$ meant to be the inverse matrices, and include also the equations 
\[
XX' = X'X = I_n \qquad YY' = Y'Y = I_m \qquad ZZ' = Z' Z = I_p,
\]
where $I_n$ denotes the $n \times n$ identity matrix, which is $\Id_U$ in any basis. (We could have instead introduced new variables such as $\delta$ and the equation $\det(X) \delta = 1$, however, the latter equation is degree $n$, whereas the above equations all have degree $O(1)$, which is more desirable from the point of view of algebraic proof complexity.)

\subsection{Polynomial encodings and the inversion principle} \label{sec:prelim:inv}
Some  principles of linear algebra can be formulated as tautologies in propositional logic and therefore also as a set of polynomial equations.
In this paper we preliminarily consider two such principles. 

\medskip

\noindent {\bf Rank Principle.} As a first example we consider a set of unsatisfiable polynomials encoding the principle that  the product of a $n\times r$ matrix $X$ by a $r\times n$ matrix $Y$ cannot be the identity matrix whenever $r<n$.  We consider variables $x_{i,k},y_{j,k}$ for $i,j\in [n]$ and $k\in [r]$, where $r< n$ to encode $X$ and $Y$. Then the polynomial encoding is:
$$\mathbb I(r,n) := \sum_{k\in[r]} x_{i,k}y_{j,k}-\delta_{i,j} \quad \quad i,j\in [n]$$
where $\delta_{i,j}=1$ if $i=j$ and $0$ otherwise. This set of polynomials is clearly unsatisfiable as long as $r<n$. 

\medskip

\noindent {\bf Inversion Principle.} The second principle encodes the invertibility of a square $n\times n$ matrix $A$,  expressing the tautology that
 $AB=I \rightarrow BA=I$ where $A,B$ are $n\times n$ matrices and $I$ is the identity matrix. Stephen A. Cook suggested this principle as a tautology that may be hard to prove in several proof systems. 
 
Let $a_{i,j},b_{i,j}$ be formal variables 
 encoding respectively  the $(i,j)$-th entries of $A$ and $B$.  We represent the fact that $AB=I$ as the set of degree $2$ polynomials

$$
\sum_{k \in [n]} a_{i,k}b_{k,j} - \delta_{i,j}\qquad i,j \in [n],
$$
where $\delta_{i,j}=1$ if $i=j$ and $0$ otherwise. We denote this set of polynomials by $AB=I$. In Section \ref{sec:matrix}, we study the degree complexity of  $AB=I \vdash BA=I$, that is of PC derivations of the polynomials $BA=I$ from the polynomials $AB=I$. 

\medskip

In view of the results we obtain in Section \ref{sec:matrix}, in Section~\ref{sec:conj} we considered a  {\em polynomial rule schema} of the form 
$$
\frac{AB=I}{BA=I}
$$
which we call the {\em Inversion Rule} (INV) meant to be added to $PC$ as an extra rule. We make this slightly more precise here.

A polynomial instantiation $\tau$ of the polynomials $AB=I$ is a substitution of polynomials $p_{i,j},q_{i,j}$ to variables $a_{i,j}$ and $b_{i,j}$.   
In PC+INV a polynomial $p$ is derivable from a set of polynomials
${\cal P}$ if 
\begin{enumerate}
\item $p$ is an axiom, or $p \in {\cal P}$;
\item $p$ is obtained by multiplication or linear combination from previous polynomials in the proof;
\item $p$ is a polynomial among a polynomial instantiation $\tau$ of $BA=I$, given that among the polynomials previously derived  in the proof there are all the polynomials 
forming the instantiation $\tau$ of $AB=I$.
\end{enumerate}

\medskip

\noindent {\bf Pigeonhole Principle.} An important role in proving the results in Section \ref{sec:matrix} is played by the well-known {\em Pigeonhole principle} stating that  any function $f$ from $[n]$ to $[r]$ with $r<n$ has a collision, that is there are $i \not= i' \in [n]$ and a $j\in[r]$ such that $f(i)=f(i')=j$. $PHP^n_r$ is the set of polynomials:

$$
\sum_{k\in [r]}p_{i,k} -1, \mbox{ for } i \in [n], \qquad \qquad p_{i,k}p_{j,k}, \mbox{ for } i \not= j \in [n], k \in [r] \qquad p_{ij}^2 - p_{ij} \text{, for $i \in [n], j \in [r]$} 
$$

Razborov \cite{DBLP:journals/cc/Razborov98} additionally included the ``functional equations'' (encoding that each pigeon cannot be matched to more than one hole):
\[
p_{i,k} p_{i,k'}, \mbox{ for } i \in [n], k \neq k' \in [r].
\]

\medskip



